% Entry point for the new two-step bridge. Existing C053 diagrams are untouched.
\begin{proofbox}
\textbf{\textcolor{CProof}{方法一的思路：第一步，先只看初中几何}}

在直角梯形 $MABN$ 中，$AM\parallel BN$。$P$ 是斜腰 $AB$ 上的普通点，
过 $P$ 作平行于底边的 $PQ$，并作垂直于底边的 $HK$。
记 $AP=r$、$PB=s$、$PQ=h$。
\begin{center}% Pure Euclidean geometry only: no conic, focus, directrix or coordinates.
% Same actual points and scale as C053_trapezoid_parabola_mapping.tex.
\begin{tikzpicture}[font=\small,>=stealth]
\draw[rounded corners=3pt,gray!45] (0,0) rectangle (14.2,11.5);
\node at (7.1,10.95) {先看一个直角梯形：上底怎样加，下底怎样减？};
\begin{scope}[shift={(1.6,5)},scale=1.15]
\fill[blue!7] (2,0)--(2,3.464101615)--(4,3.464101615)--cycle;
\fill[orange!10] (2,0)--(2,-1.154700538)--(1.333333333,-1.154700538)--cycle;
\draw[gray!80,thick] (0,3.464101615)--(0,-1.154700538);
\draw[blue!65!black,thick] (0,3.464101615)--(4,3.464101615);
\draw[orange!85!black,thick] (0,-1.154700538)--(1.333333333,-1.154700538);
\draw[blue!75!black,very thick] (2,0)--(4,3.464101615)
 node[midway,below right=5pt] {$r$};
\draw[orange!85!black,very thick] (2,0)--(1.333333333,-1.154700538)
 node[midway,above left=5pt] {$s$};
\draw[gray,dashed] (2,-1.154700538)--(2,3.464101615);
\draw[gray,dashed] (0,0)--(3.15,0);
\node[above=5pt] at (.65,0) {$h$};
\draw[blue!65!black,densely dashed] (1.333333333,-1.154700538)--(2,-1.154700538);
\draw[gray!80] (0,3.26)--(.2,3.26)--(.2,3.464101615);
\draw[gray!80] (0,-.95)--(.2,-.95)--(.2,-1.154700538);
\draw[blue!65!black] (2,3.26)--(2.2,3.26)--(2.2,3.464101615);
\draw[blue!65!black] (1.8,-1.154700538)--(1.8,-.954700538)--(2,-.954700538);
\draw[violet,->] (2.55,0) arc[start angle=0,end angle=60,radius=.55];
\node[violet] at (2.83,.34) {$\theta$};
\fill (0,3.464101615) circle(2pt) node[left=5pt] {$M$};
\fill (4,3.464101615) circle(2pt) node[above=8pt] {$A$};
\fill (1.333333333,-1.154700538) circle(2pt) node[below left=6pt] {$B$};
\fill (0,-1.154700538) circle(2pt) node[left=5pt] {$N$};
\fill (0,0) circle(1.5pt) node[left=5pt] {$Q$};
\fill (2,0) circle(2pt) node[above left=6pt] {$P$};
\fill (2,3.464101615) circle(1.5pt) node[above left=5pt] {$H$};
\fill (2,-1.154700538) circle(1.5pt) node[right=5pt] {$K$};
\node[above=6pt,blue!65!black] at (3,3.464101615) {$r\cos\theta$};
\draw[gray,dotted] (1.333333333,-1.55)--(1.333333333,-2.05);
\draw[gray,dotted] (2,-1.55)--(2,-2.05);
\draw[<->,orange!85!black] (1.333333333,-2)--(2,-2)
 node[midway,below] {$s\cos\theta$};
\end{scope}
\node[anchor=west,align=left] at (8.6,7)
 {$AM\parallel QP\parallel BN$\\[10pt]
 $AP=r,\quad PB=s$\\[9pt]$PQ=h$\\[13pt]
 $MH=NK=PQ=h$\\[10pt]
 $HA=r\cos\theta$\\[9pt]$BK=s\cos\theta$};
\node at (7.1,1.5) {看平行线和直角三角形：$AM=MH+HA$，$BN=NK-BK$。};
\node at (7.1,.6) {$\displaystyle AM=h+r\cos\theta,\qquad BN=h-s\cos\theta$};
\end{tikzpicture}
\end{center}
图示 $0<\theta<90^\circ$。在直角三角形 $PHA$、$PKB$ 中，水平直角边分别是
$r\cos\theta$、$s\cos\theta$；又 $MH=NK=PQ=h$，于是上底相加、下底相减。
\textbf{先记住这个梯形的长度分解，再给图中的点和线赋予新的含义。}
\end{proofbox}
\clearpage
\begin{proofbox}
\textbf{\textcolor{CProof}{方法一的思路：第二步，把同一个梯形放回抛物线}}

梯形的大小、形状和辅助线都与上一图相同。
现在令 $P$ 对应焦点 $F$，$MN$ 所在直线对应准线，故 $h=FQ=p$。
\begin{center}% Same trapezoid and coordinates as the pure geometry figure, translated
% from y^2=4x by x -> x+1. Here F=(2,0), directrix x=0 in drawing coordinates.
\begin{tikzpicture}[font=\small,>=stealth]
\draw[rounded corners=3pt,gray!45] (0,0) rectangle (14.2,11.5);
\node at (7.1,10.95) {再放回抛物线：只补上 $AM=AF$、$BN=BF$ 两个等量};
\begin{scope}[shift={(1.6,5)},scale=1.15]
\fill[blue!7] (2,0)--(2,3.464101615)--(4,3.464101615)--cycle;
\fill[orange!10] (2,0)--(2,-1.154700538)--(1.333333333,-1.154700538)--cycle;
\draw[teal!75!black,thick,samples=120,domain=-2.3:4.15,variable=\t]
 plot ({1+\t*\t/4},{\t});
\draw[violet,densely dashed] (0,-2.35)--(0,4.2);
\node[violet,rotate=90] at (-.65,1.5) {准线};
\draw[gray!80,thick] (0,3.464101615)--(0,-1.154700538);
\draw[blue!65!black,thick] (0,3.464101615)--(4,3.464101615);
\draw[orange!85!black,thick] (0,-1.154700538)--(1.333333333,-1.154700538);
\draw[blue!75!black,very thick] (2,0)--(4,3.464101615)
 node[midway,below right=5pt] {$r$};
\draw[orange!85!black,very thick] (2,0)--(1.333333333,-1.154700538)
 node[midway,above left=5pt] {$s$};
\draw[gray,dashed] (2,-1.154700538)--(2,3.464101615);
\draw[gray,dashed] (0,0)--(3.15,0);
\node[above=5pt] at (.65,0) {$p$};
\draw[blue!65!black,densely dashed] (1.333333333,-1.154700538)--(2,-1.154700538);
\draw[gray!80] (0,3.26)--(.2,3.26)--(.2,3.464101615);
\draw[gray!80] (0,-.95)--(.2,-.95)--(.2,-1.154700538);
\draw[blue!65!black] (2,3.26)--(2.2,3.26)--(2.2,3.464101615);
\draw[blue!65!black] (1.8,-1.154700538)--(1.8,-.954700538)--(2,-.954700538);
\draw[violet,->] (2.55,0) arc[start angle=0,end angle=60,radius=.55];
\node[violet] at (2.83,.34) {$\theta$};
\fill (0,3.464101615) circle(2pt) node[left=5pt] {$M$};
\fill (4,3.464101615) circle(2pt) node[above=8pt] {$A$};
\fill (1.333333333,-1.154700538) circle(2pt) node[below left=6pt] {$B$};
\fill (0,-1.154700538) circle(2pt) node[left=5pt] {$N$};
\fill (0,0) circle(1.5pt) node[left=5pt] {$Q$};
\fill (2,0) circle(2pt) node[above left=6pt] {$F$};
\fill (2,3.464101615) circle(1.5pt) node[above left=5pt] {$H$};
\fill (2,-1.154700538) circle(1.5pt) node[right=5pt] {$K$};
\end{scope}
\node[anchor=west,align=left] at (8.6,7)
 {上一图：$P\longrightarrow F$\\[9pt]$h\longrightarrow p$\\[14pt]
 抛物线的等距定义：\\[10pt]
 $\textcolor{blue!75!black}{AM=AF=r}$\\[10pt]
 $\textcolor{orange!85!black}{BN=BF=s}$\\[14pt]
 梯形的分解仍然成立。};
\node at (7.1,1.5) {几何关系给“加减”，抛物线定义给“相等”。};
\node at (7.1,.6) {$\displaystyle r=p+r\cos\theta,\qquad s=p-s\cos\theta$};
\end{tikzpicture}
\end{center}
抛物线定义只补上两件事：$AM=AF=r$、$BN=BF=s$。
把它们代入上一图的两条几何关系，就得到图下的两个方程。因此
\[
 r=\frac p{1-\cos\theta},\quad s=\frac p{1+\cos\theta},\qquad
 L=r+s=\frac{2p}{\sin^2\theta}.
\]
两图取 $p=2$、$\theta=60^\circ$；一般方向按有向投影写方程。
\textbf{下面保留原来的两张逐端投影图，分别对照 $AF$ 和 $BF$。}
\end{proofbox}
